\documentclass{amsart}
% For dealing with references we use the comment environment
\usepackage{verbatim}
\newenvironment{reference}{\comment}{\endcomment}
%\newenvironment{reference}{}{}
\newenvironment{slogan}{\comment}{\endcomment}
\newenvironment{history}{\comment}{\endcomment}

% For commutative diagrams we use Xy-pic
\usepackage[all]{xy}

% We use 2cell for 2-commutative diagrams.
\xyoption{2cell}
\UseAllTwocells

% We use multicol for the list of chapters between chapters
\usepackage{multicol}

% This is generally recommended for better output
\usepackage{lmodern}
\usepackage[T1]{fontenc}

% For cross-file-references

% Package for hypertext links:
\usepackage{hyperref}

% For any local file, say "hello.tex" you want to link to please

% Theorem environments.
%
\theoremstyle{plain}
\newtheorem{theorem}[subsection]{Theorem}
\newtheorem{proposition}[subsection]{Proposition}
\newtheorem{lemma}[subsection]{Lemma}

\theoremstyle{definition}
\newtheorem{definition}[subsection]{Definition}
\newtheorem{example}[subsection]{Example}
\newtheorem{exercise}[subsection]{Exercise}
\newtheorem{situation}[subsection]{Situation}

\theoremstyle{remark}
\newtheorem{remark}[subsection]{Remark}
\newtheorem{remarks}[subsection]{Remarks}

\numberwithin{equation}{subsection}

% Macros
%
\def\lim{\mathop{\mathrm{lim}}\nolimits}
\def\colim{\mathop{\mathrm{colim}}\nolimits}
\def\Spec{\mathop{\mathrm{Spec}}}
\def\Hom{\mathop{\mathrm{Hom}}\nolimits}
\def\Ext{\mathop{\mathrm{Ext}}\nolimits}
\def\SheafHom{\mathop{\mathcal{H}\!\mathit{om}}\nolimits}
\def\SheafExt{\mathop{\mathcal{E}\!\mathit{xt}}\nolimits}
\def\Sch{\mathit{Sch}}
\def\Mor{\mathop{\mathrm{Mor}}\nolimits}
\def\Ob{\mathop{\mathrm{Ob}}\nolimits}
\def\Sh{\mathop{\mathit{Sh}}\nolimits}
\def\NL{\mathop{N\!L}\nolimits}
\def\CH{\mathop{\mathrm{CH}}\nolimits}
\def\proetale{{pro\text{-}\acute{e}tale}}
\def\etale{{\acute{e}tale}}
\def\QCoh{\mathit{QCoh}}
\def\Ker{\mathop{\mathrm{Ker}}}
\def\Im{\mathop{\mathrm{Im}}}
\def\Coker{\mathop{\mathrm{Coker}}}
\def\Coim{\mathop{\mathrm{Coim}}}

% Boxtimes
%
\DeclareMathSymbol{\boxtimes}{\mathbin}{AMSa}{"02}

%
% Macros for moduli stacks/spaces
%
\def\QCohstack{\mathcal{QC}\!\mathit{oh}}
\def\Cohstack{\mathcal{C}\!\mathit{oh}}
\def\Spacesstack{\mathcal{S}\!\mathit{paces}}
\def\Quotfunctor{\mathrm{Quot}}
\def\Hilbfunctor{\mathrm{Hilb}}
\def\Curvesstack{\mathcal{C}\!\mathit{urves}}
\def\Polarizedstack{\mathcal{P}\!\mathit{olarized}}
\def\Complexesstack{\mathcal{C}\!\mathit{omplexes}}
% \Pic is the operator that assigns to X its picard group, usage \Pic(X)
% \Picardstack_{X/B} denotes the Picard stack of X over B
% \Picardfunctor_{X/B} denotes the Picard functor of X over B
\def\Pic{\mathop{\mathrm{Pic}}\nolimits}
\def\Picardstack{\mathcal{P}\!\mathit{ic}}
\def\Picardfunctor{\mathrm{Pic}}
\def\Deformationcategory{\mathcal{D}\!\mathit{ef}}

\setlength{\emergencystretch}{3em}
\begin{document}
\title[Stacks Project, Tag 070M]{319 --- K-injective limits of truncation resolutions and derived completeness}
\maketitle
\noindent Copyright (C) 2005--2025 Johan de Jong. Stacks Project authors. Source: \href{https://stacks.math.columbia.edu/tag/070M}{Tag 070M}.

\noindent Editorial difficulty note (not author text). Difficulty H2: The proof builds the ordinary inverse limit of split-surjective injective resolutions degree by degree. It constructs a termwise split product exact sequence proving that this limit is both K-injective and the derived truncation limit, then identifies the canonical comparison from the original complex. The criterion exposes precisely the obstruction to using truncation resolutions in the unbounded category..

\noindent Editorial provenance note (not author text). History: excerpt from revision \texttt{a0444\allowbreak 6e57e\allowbreak c1fbc\allowbreak 252a8\allowbreak 71afc\allowbreak ec775\allowbreak 2fb28\allowbreak 07b14}, derived.tex, lines 10588--10660. Reference presentation was adapted; original-author context and core support blocks were retained or added. The mathematical wording is unchanged. Licensed under GNU FDL 1.2 or later; no invariant sections or cover texts. See ../licenses/COPYING. Context blocks reproduce author definitions and supporting results; they are not additional counted problems.

\begin{definition}
\label{definition-K-injective}
Let $\mathcal{A}$ be an abelian category. A complex $I^\bullet$
is {\it K-injective} if for every acyclic complex $M^\bullet$ we
have $\Hom_{K(\mathcal{A})}(M^\bullet, I^\bullet) = 0$.
\end{definition}

\begin{definition}
\label{definition-derived-limit}
Let $\mathcal{D}$ be a triangulated category.
Let $(K_n, f_n)$ be an inverse system of objects of $\mathcal{D}$.
We say an object $K$ is a {\it derived limit}, or a
{\it homotopy limit} of the system $(K_n)$ if
the product $\prod K_n$ exists and there is a distinguished triangle
$$
K \to \prod K_n \to \prod K_n \to K[1]
$$
where the map $\prod K_n \to \prod K_n$ is given
by $(k_n) \mapsto (k_n - f_{n + 1}(k_{n + 1}))$. If this is the
case, then we sometimes indicate this by the notation $K = R\lim K_n$.
\end{definition}

\begin{remark}
\label{remark-map-into-derived-limit-truncations}
Let $\mathcal{A}$ be an abelian category. Let $K^\bullet$ be a
complex of $\mathcal{A}$. Then $\tau_{\geq -n}K^\bullet$ is an
inverse system of complexes and which in particular determines
an inverse system in $D(\mathcal{A})$. Let us assume that
$R\lim \tau_{\geq -n}K^\bullet$ exists. Then the canonical maps
$c_n : K^\bullet \to \tau_{\geq -n}K^\bullet$ are compatible with the
transition maps of our inverse system. By the defining distinguished
triangle of Definition \ref{definition-derived-limit} and
Lemma \href{https://stacks.math.columbia.edu/tag/0149}{lemma representable homological (Tag 0149)}
we conclude there exists a morphism
$$
c : K^\bullet \longrightarrow R\lim \tau_{\geq -n}K^\bullet
$$
in $D(\mathcal{A})$ such that the composition of $c$ with the projection
$R\lim \tau_{\geq -n}K^\bullet \to \tau_{\geq -m}K^\bullet$ is equal to $c_m$.
Now the morphism $c$ may not be unique, but we claim that whether or not
$c$ is an isomorphism is independent of the choice of $c$ (and of our choice
of the homotopy limit). Namely, for $i \in \mathbf{Z}$ and for $m > -i$
the composition
$$
H^i(K^\bullet) \xrightarrow{H^i(c)} H^i(R\lim \tau_{\geq -n}K^\bullet)
\to H^i(\tau_{\geq -m}K^\bullet) = H^i(K^\bullet)
$$
is the identity. Hence $H^i(c)$ is an isomorphism if and only if the
second map is an isomorphism. This is independent of $c$ and also
independent of the choice of the homotopy limit (as any two choices
are isomorphic).
\end{remark}

\begin{lemma}
\label{lemma-difficulty-K-injectives}
Let $\mathcal{A}$ be an abelian category with countable products and
enough injectives. Let $K^\bullet$ be a complex. Let $I_n^\bullet$ be
the inverse system of bounded below complexes of injectives produced by
Lemma \href{https://stacks.math.columbia.edu/tag/070F}{lemma special inverse system (Tag 070F)}. Then
$I^\bullet = \lim I_n^\bullet$ exists, is K-injective, represents
$R\lim \tau_{\geq -n}K^\bullet$ in $D(\mathcal{A})$,
and the following are equivalent
\begin{enumerate}
\item the map $K^\bullet \to I^\bullet$ (see proof) is a quasi-isomorphism,
\item the map $K^\bullet \to R\lim \tau_{\geq -n}K^\bullet$
of Remark \ref{remark-map-into-derived-limit-truncations}
is an isomorphism in $D(\mathcal{A})$.
\end{enumerate}
\end{lemma}

\begin{proof}
The statement of the lemma makes sense as $R\lim \tau_{\geq -n}K^\bullet$
exists by Lemma \href{https://stacks.math.columbia.edu/tag/0BK7}{lemma inverse limit bounded below (Tag 0BK7)}.
Each complex $I_n^\bullet$ is K-injective by
Lemma \href{https://stacks.math.columbia.edu/tag/070J}{lemma bounded below injectives K injective (Tag 070J)}.
Choose direct sum decompositions $I_{n + 1}^p = C_{n + 1}^p \oplus I_n^p$
for all $n \geq 1$. Set $C_1^p = I_1^p$. The complex
$I^\bullet = \lim I_n^\bullet$ exists because we can take
$I^p = \prod_{n \geq 1} C_n^p$. Fix $p \in \mathbf{Z}$.
We claim there is a split short exact sequence
$$
0 \to I^p \to \prod I_n^p \to \prod I_n^p \to 0
$$
of objects of $\mathcal{A}$. Here the first map is given by
the projection maps $I^p \to I_n^p$ and the second map
by $(x_n) \mapsto (x_n - f^p_{n + 1}(x_{n + 1}))$ where
$f^p_n : I_n^p \to I_{n - 1}^p$ are the transition maps.
The splitting comes from the map $\prod I_n^p \to \prod C_n^p = I^p$.
We obtain a termwise split short exact sequence of complexes
$$
0 \to I^\bullet \to \prod I_n^\bullet \to \prod I_n^\bullet \to 0
$$
Hence a corresponding distinguished triangle in $K(\mathcal{A})$
and $D(\mathcal{A})$. By Lemma \href{https://stacks.math.columbia.edu/tag/0BK6}{lemma product K injective (Tag 0BK6)}
the products are K-injective and represent the corresponding
products in $D(\mathcal{A})$.
It follows that $I^\bullet$ represents $R\lim I_n^\bullet$
(Definition \ref{definition-derived-limit}).
Since $R\lim I_n^\bullet \cong R\lim \tau_{\geq -n}K^\bullet$
as derived limits are defined on the level of the derived
category, we see that $I^\bullet$ represents $R\lim \tau_{\geq -n}K^\bullet$.
Moreover, the complex $I^\bullet$ is K-injective by
Lemma \href{https://stacks.math.columbia.edu/tag/090X}{lemma triangle K injective (Tag 090X)}.
By the commutative diagram of Lemma \href{https://stacks.math.columbia.edu/tag/070F}{lemma special inverse system (Tag 070F)}
and since $K^i = (\tau_{\geq -n}K^\bullet)^i$
for $n \gg 0$ we see that we get a unique map
$\gamma : K^\bullet \to I^\bullet$ such that the diagrams
$$
\xymatrix{
K^\bullet \ar[r] \ar[d]_\gamma & \tau_{\geq -n} K^\bullet \ar[d] \\
I^\bullet \ar[r] & I_n^\bullet
}
$$
commute. It follows that $\gamma$ is a map of complexes which represents
the map $c : K^\bullet \to R\lim \tau_{\geq -n}K^\bullet$ of
Remark \ref{remark-map-into-derived-limit-truncations}
in $D(\mathcal{A})$. In other words, the diagram
$$
\xymatrix{
K^\bullet \ar[r]_-c \ar[d]_\gamma & R\lim \tau_{\geq -n} K^\bullet
\ar[d]^{\cong} \\
I^\bullet \ar[r]^-{\cong} & R\lim I_n^\bullet
}
$$
is commutative in $D(\mathcal{A})$. The lemma follows.
\end{proof}
\end{document}
